\section{Related Work}
\label{sec:related}

Early jailbreak research largely focused on finding a single adversarial prompt that triggers a violating response.
In the white-box setting, attacks exploit gradient information and optimization to construct adversarial suffixes appended to harmful requests.
As model internals are unavailable for deployed systems, however, subsequent work shifted toward black-box attacks.
We organize the full space of jailbreak techniques by their generation mechanism (Table~\ref{tab:taxonomy}) and review each family below.

\begin{table}[htbp]
      \centering
      \caption{A taxonomy of jailbreak techniques.}
      \label{tab:taxonomy}
      \footnotesize
      \begin{tabular}{@{}>{\raggedright\arraybackslash}p{0.20\linewidth}>{\raggedright\arraybackslash}p{0.29\linewidth}p{0.43\linewidth}@{}}
            \toprule
            \textbf{Category} & \textbf{Representative work} & \textbf{Mechanism} \\
            \midrule
            Token / Character Manipulation    & Single Character Perturbations~\cite{scp}, CipherChat~\cite{cipherchat}                       & Perturb or encode tokens and characters (base64, unicode, added whitespace, ciphers) to attempt a jailbreak \\
            \midrule
            Gradient-Based Optimization       & GCG~\cite{gcg}, AutoPrompt~\cite{autoprompt}                                                  & Search for suffixes and triggers via gradient search to perturb internal representations (white-box)        \\
            \midrule
            Semantic Rewrite                  & PAP~\cite{pap}, HILL~\cite{hill}                                                              & Rephrase a harmful request into persuasive, natural language to evade refusal                               \\
            \midrule
            Persona Modification              & Persona Prompt~\cite{personaprompts}, Persona Modulation~\cite{personamod}, grandma exploit   & Assign the target a persona or character to induce compliance with harmful instructions                     \\
            \midrule
            Prefix Injection                  & Priming Attacks~\cite{priming}, DSN~\cite{dsn}                                                & Force the opening text of the response to bypass refusal                                                    \\
            \midrule
            Scenario Framing                  & DeepInception~\cite{deepinception}                                                            & Wrap the request in nested fictional or hypothetical scenarios to dilute its directness                     \\
            \midrule
            Context Manipulation (Multi-turn) & Crescendo~\cite{crescendo}, ASTRA~\cite{astra}, GOAT~\cite{goat}                              & Exploit accumulated dialogue context to elicit harmful responses gradually                                  \\
            \midrule
            Iterative Single-turn             & PAIR~\cite{pair}, TAP~\cite{tap}, Jailbreak-R1~\cite{jbr1}, AutoDAN-Turbo~\cite{autodanturbo} & Attacker model iteratively generates and refines prompts, delivering the best one to the target             \\
            \bottomrule
      \end{tabular}
\end{table}

\paragraph{Optimization-based attacks.}
Optimization-based attacks assume white-box access to the target model and search for adversarial inputs using gradient signals.
GCG~\cite{gcg} optimizes a transferable adversarial suffix through greedy coordinate gradient search, and AutoPrompt~\cite{autoprompt} automatically discovers trigger tokens via gradient-guided search.
Despite their high efficacy, these methods require access to the target model's parameters and gradients, which closed-source systems do not expose, together with substantial machine learning expertise.

In contrast, black-box attacks operate without access to model internals, relying solely on input-output interactions.
We review them following their progression from single-turn to multi-turn.

\paragraph{Single-turn attacks.}
Initially, black-box attacks embedded the entire malicious intent in a single, handcrafted prompt designed from human intuition.
Representative strategies include persona assignment (Persona Prompts~\cite{personaprompts}, Persona Modulation~\cite{personamod}, PAP~\cite{pap}), nested fictional or hypothetical framing (DeepInception~\cite{deepinception}), persuasive semantic rewriting (HILL~\cite{hill}), prefix injection (Priming attacks~\cite{priming}, DSN~\cite{dsn}), and token-level or encoding-level manipulation (Single Character Perturbations~\cite{scp}, CipherChat~\cite{cipherchat}).
Since the entire intent is packed into one static prompt, these attacks exhibit limited diversity and are easily detected and patched.

\paragraph{Iterative single-turn attacks.}
To reduce manual effort and improve robustness, later work automated prompt construction using an attacker LLM that iteratively generates and refines candidate prompts based on the target's feedback (PAIR~\cite{pair}, TAP~\cite{tap}, Jailbreak-R1~\cite{jbr1}).
These approaches still deliver a single-turn prompt and do not exploit the conversational context that accumulates across turns.

\paragraph{Multi-turn attacks.}
In practice, an adversary rarely succeeds in a single attempt; instead, the attack unfolds gradually over a conversation.
Multi-turn attacks exploit this by distributing the malicious intent across several turns, keeping each turn benign while steering the model toward a harmful goal.
Building on this idea, recent work develops automated attack agents: Crescendo~\cite{crescendo} escalates from benign queries, GOAT~\cite{goat} casts multi-turn red teaming as a generative offensive agent, and ASTRA~\cite{astra} discovers, retrieves, and evolves attack strategies.
A more recent line builds adaptive multi-turn agents that plan and refine an attack trajectory: ActorAttack~\cite{actorattack} derives attack chains from an actor-network built around the goal, MRJ-Agent~\cite{mrjagent} decomposes a harmful request across rounds under risk control, and X-Teaming~\cite{xteaming} coordinates planner, attacker, and verifier agents to search multi-turn scenarios.
Several frameworks further accumulate a persistent strategy memory to refine attacks over time (AutoDAN-Turbo~\cite{autodanturbo}, AutoRedTeamer~\cite{autoredteamer}), while Rainbow Teaming~\cite{rainbow} maximizes prompt diversity through a quality-diversity archive of single-turn adversarial prompts.
Each addresses only part of the problem. Rainbow optimizes diversity but stays single-turn and target-agnostic, whereas the memory frameworks build one strategy library shared across targets. AutoDAN-Turbo, for instance, retrieves by response similarity rather than by the model under attack, so the same pool is applied regardless of target.
These adaptive agents refine an attack within a single conversation, but none carries a per-target success memory across goals or ranks strategies by a per-target success estimate, so no prior method concentrates its attacks on what has already worked against a specific model. This leaves prompt diversity and cross-goal target adaptation jointly unresolved, the gap our work closes.
